\section{\faUsers\ 工作经历}

\datedsubsection{WeRide 感知组研发工程师}{2024.06 - 至今}

\paragraph{AEB 障碍物感知融合}
\begin{itemize}
      \item 参与奇瑞 \textit{single Orin Y} 车型和广汽 \textit{G300} 车型的 AEB 功能开发，设计融合算法。
      \item ASIL-B 要求通过输入冗余实现容错。
            设计多来源融合算法，融合视觉端到端模型、视觉 AEB 模型和 Radar 检测（\textit{G300}）。
      \item 基于 Kalman Filter 融合多来源的检测。通过量测噪声解耦各个输入的差异。通过预测对齐时间戳。
      \item Orin Y 车型交付达到100\% UT 覆盖，仿真 5kh 无 fp，通过 C-NCAP 考试。
\end{itemize}

\paragraph{LiDAR freespace 后处理}
\begin{itemize}
      \item 雨天场景下高反牌（路牌、锥桶等）会导致 LiDAR 感知产生强烈误检。
            设计基于规则的后处理来移除误检。
            通过 PCA 拟合高反牌点云所在平面，移除由高反牌造成的 freespace 误检。
            并结合视觉模型输出进行补偿。
      \item 实现基于 grid tracking 的 freespace记忆模块。
            通过盲区穿越判定策略补全车后可行驶区域，
            相比 Log-Odds 累计概率结果稳定且可预测。
            用于支持三点掉头等复杂场景；能够过滤偶发的误检。
      \item 设计用于车端HMI渲染的 freespace 点云到 linestring 的转换算法。
            基于fs记忆模块提供的稳定结果，
            通过不同的密度对角度进行量化，采样结果求凸包，最后在凸包顶点间补充 fs 点丰富细节。
            使车端HMI能够渲染 fs 轮廓线，并兼顾一定精度。
\end{itemize}

\paragraph{NeuralMap 车道线后处理}
\begin{itemize}
      \item NeuralMap 是 WeRide compass 无图项目的车道线感知模块。
            用于与奇瑞、博世等合作的所有 L2 无图车型。
      \item 根据下游反馈和需求开发并调整后处理功能。
            赋值分合流、车道线类型等属性。
      \item 基于  Kalman Filter 识别并平滑 crosswalk，为 HMI 提供稳定的结果。
      \item 开发基于bag的自动化标注工具,
            实现自车行为(车道居中/偏离、变道、停止线前静止/起步/穿越)
            和交互车辆行为(静止、加减速、切入、绕行)的自动识别和标注。
\end{itemize}



\datedsubsection{TuSimple 研发工程师——高精地图 full time \& intern }{2022.07 - 2024.03}

\paragraph{高精地图 SDK TSMap 的维护和开发}
\begin{itemize}
      \item TSMap 是 TuSimple 的高精地图 C++ SDK。同时通过 pybind11 提供 Python 接口。
      \item 开发 Perception Drivable Map。
            提供接口给感知快速获取表达可行驶区域的布尔矩阵。
      \item 开发 Planning Drivable Map。
            计算 freespace 距离。利用局部性异步计算结果以满足实时性要求。
      \item 根据下游场景进行一些优化。例如提供批量查询曲率接口，优化了批量查询的复杂度。
\end{itemize}

\paragraph{地图产线}
\begin{itemize}
      \item 参与用于地图量产的微服务架构的新地图产线开发。
      \item 完成大部分车道线构建相关服务的开发。
            流程输入车道线矢量数据，输出包含车道、路口、拓扑关系、物理元素关系等地图元素的母库规格数据。
            接入可观测性平台，便于快速发现和溯源问题。
      \item 维护旧地图编译工具。
            根据算法和地图需求，拓展了地图元素类型，支持表达双线型边界、非车道可行驶区域等，且保证了向前和向后兼容。
\end{itemize}


